\documentclass[11pt,a4paper]{article}
\usepackage[T2A]{fontenc}
\usepackage[utf8]{inputenc}
\usepackage[russian,english]{babel}
\usepackage{amsmath,amssymb,amsthm,mathtools}
\usepackage{setspace}
\usepackage{enumitem}
\usepackage[hidelinks]{hyperref}
\setstretch{1.3}
\hyphenpenalty=10000
\exhyphenpenalty=10000
\renewcommand{\leq}{\leqslant}
\renewcommand{\le}{\leqslant}
\renewcommand{\geq}{\geqslant}
\renewcommand{\ge}{\geqslant}
\newcommand{\head}[1]{\par\bigskip\noindent\textbf{\underline{#1}}\quad}
\begin{document}
\begin{center}
{\huge\bfseries Limits -- 9\par}\vspace{1.2em}
{\Large Sequences defined by integrals and sums\par}\vspace{0.8em}
{\foreignlanguage{russian}{1 сентября 2026 года}\par}
\end{center}

\noindent\textbf{\underline{Theoretical minimum.}}\par\vspace{-\bigskipamount}\medskip

\head{Theorem 1 (Riemann sums).} If $f$ is integrable on $[0,1]$, then
$\frac1n\sum_{k=1}^nf\bigl(\frac kn\bigr)\to\int_0^1f$; for monotone $f$ the error is at most
$\frac{|f(1)-f(0)|}{n}$, and for smooth $f$ it equals $\frac{f(1)-f(0)}{2n}+O(n^{-2})$.

\head{Theorem 2 (oscillation and averaging).} If $g$ is periodic with period $T$ and integrable,
then $\int_a^bf(x)g(nx)\,dx\to\frac1T\int_0^Tg\cdot\int_a^bf$; in particular
$\int_a^bf(x)\sin(nx)\,dx\to0$ (the Riemann--Lebesgue lemma).

\head{Fact 1.} If the mass of the integral concentrates near the point of maximum, Laplace's method
applies: $\int_0^1f(x)^ndx\sim C\,\frac{f(x_0)^n}{\sqrt n}$, where $x_0$ is the point of maximum;
the simplest instance is $\int_0^{\pi/2}\sin^nx\,dx\sim\sqrt{\frac\pi{2n}}$.

\noindent\rule{\textwidth}{0.4pt}\par\medskip

\begin{enumerate}[itemsep=0.6em]
\item Prove Theorem 1 together with the error estimate, and compute
$\lim_{n\to\infty}\frac1n\sum_{k=1}^n\frac1{1+(k/n)^2}$ and $\lim_{n\to\infty}\sum_{k=1}^n\frac n{n^2+k^2}$.

\item Prove Theorem 2 first for a step function $f$ and then for an integrable one; deduce the
Riemann--Lebesgue lemma and compute $\lim_{n\to\infty}\int_0^\pi\frac{\sin x}{1+3\cos^2(2nx)}\,dx$.

\item Let $a_n=\int_0^1\frac{x^n}{1+x}\,dx$. Find the asymptotics of $a_n$, prove that
$\sum_{n\ge1}(-1)^na_n$ converges and compute its sum.

\item (IMC 2025 shortlist) Let $f\colon[0,+\infty)\to[0,+\infty)$ be continuous with $f(x)\to C>0$
as $x\to+\infty$. For every positive integer $n$ put $a_n=\int_{n^2}^{n^2+1}\bigl(f(t)\bigr)^ndt$ and
$b_n=\int_{(n-1)^2}^{(n-1)^2+1}\bigl(f(t)\bigr)^ndt$. Suppose that $a_1+a_2+\dots+a_n\to\infty$ and
$\frac{a_{n+1}}{a_1+a_2+\dots+a_n}\to L$. Find
$\lim_{n\to\infty}\frac{b_1+b_2+\dots+b_n}{a_1+a_2+\dots+a_n}$ in terms of $C$ and $L$.

\item Prove by Laplace's method that $\int_0^{\pi/2}\sin^nx\,dx\sim\sqrt{\frac\pi{2n}}$ and deduce
Wallis' formula; compare with the exact answer through the gamma function.
\end{enumerate}
\end{document}
